\documentclass[runningheads]{llncs}

\usepackage[T1]{fontenc}
\usepackage[utf8]{inputenc}
\usepackage{booktabs}
\usepackage{amsmath}
\usepackage{amssymb}
\usepackage{listings}
\usepackage{xcolor}
\usepackage{hyperref}
\usepackage[nolist]{acronym}

\lstset{
  basicstyle=\ttfamily\small,
  columns=fullflexible,
  keepspaces=true,
  mathescape=true,
  keywordstyle=\bfseries,
  morekeywords={trace,for,do,check,invs,current,counter,stores,maps,yield,int,while,if},
  literate={⌊}{{$\lfloor$}}1 {⌋}{{$\rfloor$}}1 {≤}{{$\le$}}1 {⇒}{{$\Rightarrow$}}1,
  xleftmargin=1em,
  aboveskip=0.5ex,
  belowskip=0.5ex,
}

\newcommand{\cuvee}{Cuvée}
\newcommand{\lf}[1]{\lfloor #1 \rfloor}
\newcommand{\code}[1]{\texttt{#1}}
\newcommand{\todo}[1]{\textcolor{red}{[TODO: #1]}}

\acrodef{TEQ}{Trace Equation System}
\acrodefplural{TEQ}[TEQs]{Trace Equation Systems}

\begin{document}

\title{Cuvée: Loop Summaries by Trace Equation Systems\\ and a Lemma Database\\ (Competition Preview)}
\titlerunning{Cuvée (Competition Preview)}

\author{Gidon Ernst}
\authorrunning{G. Ernst}
\institute{LMU Munich, Germany\\
\email{gidon.ernst@gmail.com}}

\maketitle

\begin{abstract}
\cuvee{} proves the safety of C programs with loops without inferring loop invariants.
Each loop is summarized by a \ac{TEQ}, which describes the sequences of values that
the variables take in the iterations. A database of lemmas, most of them proved in Lean, rewrites these equations into closed forms, such as counters, sums, and arrays filled with values.
Goals that remain after rewriting are encoded in first-order logic and passed to Z3.
\cuvee{} does not search for counterexamples: it reports a violation only when rewriting refutes a goal.
In our experiments, it proves 113 safe tasks of ReachSafety-Arrays
that none of five SV-COMP~2026 participants proves, and without the fallback to Z3 it reports no wrong
results.
\end{abstract}
\acresetall

\section{Verification Approach}

The difficult part of the verification of loops is to find inductive invariants. Invariants describe the
state at the loop head. However, the reasons why a program is correct are often easier to state about the
whole computation of the loop: \emph{the array is a sequence of stores at the positions $0, \dots, n-1$}, or
\emph{the variable is the sum of the elements read so far}. \cuvee{} follows this view. It represents the
computation of each loop by \acp{TEQ} and proves goals by rewriting with a database of lemmas
about such traces. The approach is described in detail in~\cite{teq}; we summarize it here.

\paragraph{Trace Equation Systems.}
\cuvee{} symbolically executes each function, and inlines the calls of functions without a contract. A loop that runs $n$ times is
replaced by a system of equations over \emph{traces}, sequences of values:
\begin{itemize}
\item a synthetic counter $K = \code{counter}(0, 1, n)$ of the iterations;
\item for each modified variable $x$, an equation
  \[X = \code{trace}(x_0, \code{for } \lf{K}, \lf{Y}, \dots \code{ do } x, k, y, \dots \Rightarrow u).\]
  The update $u$ has one guarded case per path through the body, and $\lf{Y}$ is the trace of another
  variable without its last value;
\item the iteration constraint $\code{invs}(\code{for } \dots \code{ check } \dots \Rightarrow
  \varphi)$, which states that each iteration ran with a true loop test along one of these paths.
\end{itemize}
The variable then has the value $\code{current}(X)$ after the loop. Nondeterministic values become fresh
traces. The system is sound and relatively complete: it describes exactly the reachable states after the
loop~\cite{teq}. Figure~\ref{fig:example} shows the \ac{TEQ} of a loop that fills an array.

\begin{figure}[t]
\begin{minipage}[t]{0.31\textwidth}
\begin{lstlisting}[language=C,basicstyle=\ttfamily\footnotesize]
int i = 0;
while (i < N) {
  a[i] = 0;
  i = i + 1;
}
// assert a[j] == 0
// for 0 $\le$ j < N
\end{lstlisting}
\end{minipage}%
\begin{minipage}[t]{0.69\textwidth}
\begin{lstlisting}[basicstyle=\ttfamily\footnotesize]
K == counter(0, 1, n)
I == trace(0, for ⌊K⌋ do i, k :: i + 1)
A == trace(a, for ⌊K⌋, ⌊I⌋ do a,k,i :: a[i := 0])
invs(for ⌊K⌋, ⌊I⌋ check k, i :: i < N)
!(current(I) < N)
\end{lstlisting}
\end{minipage}
\caption{A loop (left) and its \ac{TEQ} (right), simplified. The goal is
$\code{current}(A)[j] = 0$.}
\label{fig:example}
\end{figure}

\paragraph{Rewriting with lemmas.}
\cuvee{} does not compute fixpoints or interpolants. Instead, its simplifier rewrites the equations with
lemmas of two kinds.
\begin{itemize}
\item \emph{Introduction rules} recognize a fragment of a computation and name it. For example,
  \code{counter\_def} turns the trace $I$ into $\code{counter}(0, 1, n)$, and \code{stores} turns $A$
  into the stores of zeros at the positions $\lf{I}$.
\item \emph{Acceleration rules} evaluate observers of these functions. For example, the last value of
  a counter is $\code{current}(\code{counter}(i, c, n)) = i + n \cdot c$, and the iteration constraint
  $\code{invs}(\code{for } \lf{\code{counter}(0,1,n)} \code{ check } k \Rightarrow k < N)$ is
  equivalent to $n \le N$ for $N \ge 0$. The rule \code{current\_stores} reads the value stored last at
  a position.
\end{itemize}
In the example, eleven rule applications close the goal.
The database contains 205 lemmas in eight modules, over the definitions of a core module: arithmetic,
counters, aggregation (sums, counts, minima, maxima), arrays, search, nested loops, sorting, and strided
stores. The order of the modules is the order in which rules are tried; \code{-teq} loads the first
five, and the others are loaded on demand (\code{-include teq:nested}). 164 lemmas, those of arithmetic,
aggregation, arrays, and search and most of those of counters, have machine-checked proofs in
Lean~4~\cite{lean4}. The matcher works modulo the binding structure of the
traces, modulo associativity and commutativity of $+$, and with higher-order patterns in Miller's fragment.
It also uses the context of each update, such as $0 \le k < n$ and the iteration constraints. Thus a lemma
applies to all loops that share its pattern, regardless of the names, order, and number of the other
variables.

\paragraph{Fallback to SMT.}
The remaining goals are lowered into first-order logic, where each trace is a function of the
iteration index and each equation of a \ac{TEQ} a quantified characterization of its steps, and passed to
Z3~\cite{z3}. Functions of the database without a first-order definition are generalized, so a model
does not refute the program: the verdict is then \emph{unknown}.

\section{Software Architecture}

\cuvee{} is written in Scala~3 (about 22{,}000 lines). It is a pipeline of sources, transformations,
and sinks over a common core of logic and programs, with frontends for SMT-LIB, SV-LIB, Boogie, Dafny,
C, and SV-COMP task definitions. The C frontend handles integer programs with arrays, and reads ACSL
contracts and loop invariants. Signed integers are mathematical integers, since their overflow is
undefined, whereas unsigned arithmetic wraps around modulo $2^k$; nondeterministic values lie within
the range of their type. Arrays are values, also those allocated by \code{malloc} with a symbolic size.
The transformation \code{-teq} generates the goals and loads the lemma database, which is written in
Dafny syntax. The simplifier normalizes arithmetic into polynomials, simplifies formulas under their
context, and applies the lemmas. The context decides difference constraints $x - y \le c$ among its
facts, closed under shortest paths by an incremental Floyd--Warshall algorithm, whose negative cycles
make the context inconsistent; it also bounds remainders modulo $2^k$ of the unsigned arithmetic.
\code{-split} separates the goals, and \code{-fallback} lowers the remaining ones for Z3~4.16.

\section{Strengths and Weaknesses}

We ran \cuvee{} and five participants of SV-COMP~2026, in their competition versions and
configurations, on ReachSafety-Arrays (440 tasks, 320 safe) and ReachSafety-Loops (764 tasks, 555 safe),
with BenchExec~\cite{benchexec}, 900\,s of CPU time, 7\,GB of memory, and 2 cores per run
(Table~\ref{tab:results}). \cuvee{} ran with \code{-verify} and all modules of the lemma database. The
results of the other tools are close to the official ones. We did not validate witnesses.

\begin{table}[t]
\centering
\caption{Results of \cuvee{} and five SV-COMP~2026 participants in our runs. ``Unique'' counts
the safe tasks that no other tool in the table proves. The score follows the SV-COMP scheme
($+2$ / $+1$ for correct true / false, $-32$ / $-16$ for wrong true / false), without witness validation.}
\label{tab:results}
\setlength{\tabcolsep}{3pt}
\begin{tabular}{lrrrrrrr}
\toprule
& \cuvee{} & \cuvee{}+Z3 & Symbiotic & AISE & UAutom. & CPAchecker & Korn \\
\midrule
\multicolumn{8}{l}{\emph{ReachSafety-Arrays}} \\
\quad correct true  & 137 & \textbf{147} & 68 & 65 & 11 & 4 & 34 \\
\quad correct false & 9 & 51 & 86 & \textbf{89} & 76 & 70 & \textbf{89} \\
\quad wrong & 0 & 1 & 0 & 8 & 0 & 0 & 0 \\
\quad unique true & \multicolumn{2}{c}{113} & & & & & \\
\quad score & 283 & \textbf{329} & 222 & 75 & 98 & 78 & 157 \\
\midrule
\multicolumn{8}{l}{\emph{ReachSafety-Loops}} \\
\quad correct true  & 126 & 280 & 418 & \textbf{420} & 378 & 348 & 388 \\
\quad correct false & 16 & 28 & 169 & 166 & 130 & 133 & \textbf{171} \\
\quad wrong & 0 & 8 & 0 & 0 & 0 & 1 & 1 \\
\quad unique true & \multicolumn{2}{c}{3} & & & & & \\
\quad score & 268 & 348 & 1005 & \textbf{1006} & 886 & 813 & 915 \\
\bottomrule
\end{tabular}
\end{table}

\paragraph{Strengths.}
\begin{itemize}
\item \emph{Arrays of unbounded size.} \cuvee{} proves more safe tasks of ReachSafety-Arrays than any other
  tool in the table: 147 with Z3 and 137 without, compared with 68 for the next tool. Of these, 113 are not
  proved by any of the other five tools. Most of them are in \code{array-fpi} (56), \code{array-examples}
  (25), and \code{array-crafted} (18), whose loops fill arrays with values that depend on the position or
  on earlier values, and whose sizes are symbolic. With \cuvee{}, the six tools prove 209 of the
  320 safe tasks, instead of 96. The summaries do not depend on the number of iterations, so large or
  symbolic bounds do not matter.
\item \emph{Few wrong results.} Without the fallback, \cuvee{} reports no wrong result on all 1204 tasks.
  Most lemmas are proved in Lean~4, so the trusted base is mainly the translation into
  \acp{TEQ}, the simplifier, the matcher, and the lemmas without a proof. The only wrong results,
  9 of the fallback, are discussed below.
\item \emph{Explanations.} With \code{-rules}, \cuvee{} reports the lemmas that each proof uses, and
  a residual goal points to the lemma that is missing.
\item \emph{Incremental.} Rewriting also simplifies goals that it does not close. The remaining goals are
  small enough for Z3: the fallback adds 154 proofs in Loops, 122 of them in the nonlinear tasks of
  \code{nla-digbench} and its scaled variants, where the simplifier eliminates the loops and Z3 solves the
  remaining arithmetic.
\end{itemize}

\paragraph{Weaknesses and limitations.}
\begin{itemize}
\item \emph{Few violations.} \cuvee{} does not search for counterexamples. It reports \emph{false} only
  if rewriting refutes a goal, without an abstraction on the way; a model of Z3 for a lowered goal is no
  counterexample, since the lowering generalizes. Thus it reports only 25 of the 329 unsafe tasks of both
  categories, and 79 with Z3. Nor does it produce witnesses yet. In the official scoring, its proofs
  would therefore count as unconfirmed.
\item \emph{Wrong results of the fallback.} The 7 wrong proofs are the variants of
  \code{nla-digbench-scaling/fermat1-ll\_unwindbound}, whose shared counter exits three nested loops
  early; presumably the lowering for Z3 loses a constraint of these exits. The 2 wrong violations come
  from goals that the fallback lowers, whose refutation by Z3 must not count as a counterexample.
\item \emph{Machine integers.} Signed integers are mathematical integers, since their overflow is
  undefined behavior, which the tasks are meant to avoid. The wrap-around of unsigned integers is
  modeled by remainders modulo $2^k$, which make goals harder to close.
\item \emph{Coverage of the database.} A loop is summarized only if the database has a lemma for its shape.
  Loops whose summary needs only partial facts, such as bounds, minima, or maxima under conditional
  updates, are often not closed. In Loops, many tasks need nonlinear invariants, such as
  \code{nla-digbench}, or their loops are state machines. Here \cuvee{} without Z3 proves only 126 tasks,
  while the other tools prove 348 to 420. Each lemma fixes the shape of its left-hand side, so a loop that
  writes $s$ cells under $s$ independent conditions requires a lemma for each $s$. Rewriting may also
  diverge: 58 tasks end with \emph{rewriting does not terminate}.
\item \emph{Frontend.} \cuvee{} rejects floating-point numbers (29 tasks), \code{goto} into code that
  falls through (26), and fails to parse 37 tasks, mostly because of structs and pointers other than
  arrays. Bitwise operations are uninterpreted functions, except shifts by constants.
\item \emph{Limits of the fallback.} The first-order encoding does not cover every goal, e.g., when a goal
  has sequences that are not bound by a quantifier (8 tasks) or lists nested more than twice (8). More
  often, Z3 does not answer in time: the fallback times out on 336 tasks, against 17 without it.
\item \emph{Speed.} The median CPU time of a proof on Arrays is 20\,s, with or without Z3, compared with
  4\,s for Symbiotic. Even a trivial program takes about 5\,s of CPU time, for the start of the
  JVM and the loading of the lemma database.
\end{itemize}

Many lemmas of the database, among them those for conditional stores, sums with resets, and strided
stores, were found on the MonoCera benchmarks~\cite{monocera,monocera-artifact}.

\section{Tool Setup and Configuration}

\cuvee{} needs Java~21. The archive contains \code{cuvee.jar}, the script \code{Cuvee.sh}, and a Z3
binary, which the script puts on the path. The BenchExec tool-info module is \code{cuvee.py} (\code{benchexec.tools.cuvee}). It runs
\begin{center}
\small\code{./Cuvee.sh -teq -simplify -split -simplify -fallback -z3}\\
\small\code{-timeout 900000 -verify [-arch 32|64] task.c}
\end{center}
where \code{-arch} follows the data model of the task. With \code{-verify}, \cuvee{} answers
\emph{correct} if all goals are proved, \emph{incorrect} if a goal is refuted without abstraction, and
\emph{unknown} otherwise, which \code{cuvee.py} reports as \emph{true}, \emph{false}, and \emph{unknown}.
For C, it supports only the property \code{unreach-call}, and it participates in ReachSafety-Arrays and
ReachSafety-Loops. For SV-LIB scripts, \code{-verify} answers each \code{verify-call}.

\cuvee{} also participates in the SV-LIB categories CoreVerification and CTranslated of the property
\code{correct-tags}, with the same command on the script. The tool-info module reports \emph{false}
if a \code{verify-call} is answered \emph{incorrect}, and \emph{true} if all of them are answered
\emph{correct}. The archive also contains the licenses and a smoke test, which runs \code{-teq}, with
and without the fallback to Z3, and \code{-chc -z3}, the encoding into constrained Horn clauses, on a
safe and an unsafe C program and SV-LIB script each.

\section{Software Project and Contributors}

\cuvee{} is developed by Gidon Ernst at LMU Munich, at
\url{https://github.com/gernst/cuvee-plus}, under the Apache License~2.0. It is built with Mill.
The lemma database is in the folder \code{resources/cuvee/teq} of the repository, and the
Lean proofs are in the file \code{resources/cuvee/teq.lean}.

\paragraph{Data-Availability Statement.}
The version of \cuvee{} used in SV-COMP~2027 is available in the SV-COMP tool archive~\cite{cuvee-artifact}.

\section*{AI Declaration}

The theory of \acp{TEQ}, the translation of loops, the rewriting strategy, the matching of
traces, and the soundness argument, was fully worked out before this implementation~\cite{teq}. An
earlier implementation in Scala~2, with a lemma database of about 65 lemmas written in Boogie, produced
the results of that work. The present version of \cuvee{} started from a Scala~3 code base of about
3{,}500 lines written by the author: the core of logic and programs, the pipeline, and the frontends for
SMT-LIB and Dafny.

From there, coding agents (Claude Code with Claude Opus) carried out most of the implementation, under
the direction of the author, who set the design, decided the open questions, and reviewed the results;
574 of the 584 commits since then are co-authored by the agent. The agents wrote the generation of
\acp{TEQ}, the simplifier with its reasoning about the context, the rewriting with the database
and the matching of traces, the fallback to Z3, the encoding into constrained Horn clauses, the frontends
for Boogie, C with ACSL, SV-LIB, and SV-COMP task definitions, the tests, and the scripts of the
experiments. The code base grew to about 22{,}000 lines.

For the lemma database, the agents first ported the lemmas of the Boogie implementation to Dafny syntax,
added missing preconditions on bounds, and corrected some of them. They then extended the database to 205
lemmas in eight modules, mostly by analyzing the goals left open on the SV-COMP and MonoCera benchmarks
and writing lemmas for the missing loop shapes, under the review of the author. They also wrote the
formalization of the core definitions in Lean~4 and the proofs of 164 lemmas; the others are trusted.
The agents ran the experiments and drafted parts of this paper under the direction of the author.

\bibliographystyle{splncs04}
\bibliography{references}

\end{document}
